\section{Резюме}

Настоящата дипломна работа е посветена на проектирането и разработването на уеб базирана система за електронна търговия с компютърна и офис техника. Темата е избрана поради високата практическа значимост на онлайн търговските платформи и поради факта, че подобен тип системи съчетават в себе си съществени проблеми на съвременното софтуерно инженерство: многослойна архитектура, моделиране на предметна област, обработка на HTTP заявки, изграждане на потребителски интерфейс, валидация на входните данни, работа с релационна база данни, управление на състояние и интеграция с външни услуги. Изследваната и реализирана система има учебно-приложен характер, но покрива пълен набор от реалистични функционалности, характерни за малък до среден електронен магазин.

Практическата реализация е разработена с езика Java 11 и рамката Spring Boot 2.6.2. Обработката на уеб заявките и връзката между бизнес логиката и потребителския интерфейс са реализирани чрез Spring MVC, а визуалният слой е изграден със сървърно рендирани Thymeleaf шаблони. За достъп до данните се използват Spring Data JPA и MySQL, като моделът на приложението е представен чрез entity класове, repository интерфейси и service слой. За допълнително повишаване на качеството на разработката са приложени Hibernate Validator за входна валидация, ModelMapper за преобразуване между различни моделни представяния и Spring Security Crypto за защитено хеширане на пароли чрез \texttt{Pbkdf2PasswordEncoder}. Изпращането на имейл за потвърждение на регистрацията е реализирано посредством инфраструктурата на Spring за работа с електронна поща.

Разработената система поддържа регистрация и вход на потребители, потвърждаване на акаунт чрез електронна поща, редакция на профилни данни, управление на фирми, добавяне и категоризация на продукти, разглеждане на продуктов каталог, добавяне на артикули в количка, оформяне на поръчка, съхраняване на история на покупките и извличане на примерна агрегирана информация чрез REST крайна точка. По този начин проектът демонстрира завършен потребителски поток, започващ от създаване на акаунт и завършващ с реализиране и архивиране на поръчка.

В документа са разгледани последователно контекстът на задачата, сравнението със съществуващи решения, теоретичните постановки, изискванията към системата, архитектурната организация, моделът на данните, конкретните механизми на реализация, потребителските сценарии, конфигурацията на средата за изпълнение и подходът към верификация и оценка. Особено внимание е отделено на връзката между структурата на сорс кода в директорията \texttt{src}, генерираните артефакти в \texttt{target}, и логическите отговорности на отделните слоеве в приложението. Наред с положителните характеристики на решението са анализирани и реалните ограничения на текущата имплементация, включително използването на сесийно управление без пълноценна роля-базирана авторизация, наличие на чувствителна конфигурация в локални файлове и необходимост от по-строга изолация на част от бизнес операциите на ниво потребител.

Основният принос на дипломната работа е в систематичното описание на реален Java Spring Boot проект като цялостна инженерна разработка: от предметната област и технологичния избор, през архитектурния модел и модела на данните, до практическите потребителски сценарии и препоръките за бъдещо развитие. Полученият резултат представлява работеща основа за електронен магазин и едновременно с това илюстрира добри практики при структуриране на малка уеб информационна система. Документацията има за цел не само да представи крайния резултат, но и да обоснове защо системата е организирана по този начин, какви проблеми решава и как може да бъде развита в по-зряла посока.

\subsection{Abstract}

This thesis presents the design and development of a web-based e-commerce system for computer and office equipment. The project is motivated by the practical importance of online commerce platforms and by the fact that such systems combine several key software engineering concerns, including layered architecture, domain modeling, HTTP request processing, user interface generation, input validation, relational data persistence, session management, and integration with external infrastructure services.

The implemented application is built with Java 11 and Spring Boot 2.6.2. Spring MVC is used for request handling, Thymeleaf for server-side rendered views, Spring Data JPA and MySQL for persistence, Hibernate Validator for input constraints, ModelMapper for model transformation, and Spring Security Crypto for password hashing. The system supports user registration and login, e-mail confirmation, profile management, supplier company management, product addition and categorization, product catalog browsing, shopping cart operations, checkout, order history, and a simple REST endpoint for retrieving the maximum product price.

The thesis documents the system from multiple perspectives: problem context, comparison with existing solutions, theoretical foundation, requirements, architecture, data model, implementation, user interaction, build configuration, verification, and future work. A major emphasis is placed on mapping the actual project structure and source code organization to the logical responsibilities of the application layers. The work also identifies current limitations, such as session-based access control, hard-coded local configuration values, and opportunities for improving robustness, security, and extensibility.

The main contribution of the thesis lies in the systematic documentation of a real Spring Boot application as a complete engineering solution. The resulting system provides a functional educational prototype of an online store and a solid basis for future expansion toward a more secure and production-ready platform.

\subsection{Ключови думи}

Електронна търговия, уеб приложение, Spring Boot, Spring MVC, Thymeleaf, Spring Data JPA, MySQL, онлайн магазин, многослойна архитектура, Java.
